\documentclass[11pt,letterpaper]{article}
\usepackage{amsmath,amssymb}
\usepackage{geometry}
\usepackage{fancyhdr}
\usepackage{xcolor}
\usepackage{listings}
\usepackage{array}
\geometry{margin=1.1in}
\definecolor{certgreen}{RGB}{0,130,60}
\definecolor{codebg}{RGB}{245,245,245}
\lstset{basicstyle=\small\ttfamily,breaklines=true,frame=single,
        backgroundcolor=\color{codebg},xleftmargin=1em,aboveskip=4pt,belowskip=4pt}
\pagestyle{fancy}\fancyhf{}
\fancyhead[L]{\small\textbf{Module 3: Q5/P5 Bound} \texttt{[Q5P5Bound]}}
\fancyhead[R]{\small Machine Certificate v1.6}
\fancyfoot[C]{\thepage}
\renewcommand{\headrulewidth}{0.4pt}
\begin{document}

\begin{center}
{\Large\bfseries Module 3: Continued Fraction Obstruction for $\pi/10$}\\[4pt]
{\large Certifies: $a_6 = 733$, $Q_5 = 226$, $a_7 = 11$, bound $= 82{,}829$}\\[2pt]
{\normalsize Machine Certificate v1.6 \quad David Fox \quad May 21, 2026}
\end{center}
\vspace{0.5em}
\begin{center}
\colorbox{certgreen!15}{\parbox{0.85\textwidth}{\centering
\textcolor{certgreen}{\bfseries STATUS: CERTIFIED \checkmark}\\
Source and stdout SHA-256 verified.
}}
\end{center}

\vspace{0.5em}
\noindent\textbf{Symbol disambiguation (Battle Plan v1.6):}
$Q_n = $ denominator $q_n$ of the $n$th convergent of $\pi/10$ (NOT $6^4 = 1296$).
$a_n = $ $n$th partial quotient of $\pi/10$.
Values $Q_5 = 1296$ and $a_7 = 1$ are deprecated and shall not appear.

\section*{1. Claim}

The simple continued fraction of $\pi/10$ is
\[
  \pi/10 = [0;\; 3,\, 5,\, 2,\, 5,\, 1,\, 733,\, 11,\, 1,\, 1,\, \ldots]
\]
with convergents $p_n/q_n$.  In particular, $a_6 = 733$, $q_5 = 226$, $a_7 = 11$.
By Lemma 3.2 of Paper 6, for $n \geq 5$ the numerators satisfy
$p_n > a_{n+1} \cdot q_n / 2$.  Hence for $n = 5$:
\[
  p_5 > 733 \times 226 / 2 = 82{,}829.
\]

\section*{2. Certified Values}

\begin{center}
\begin{tabular}{|l|l|}
\hline
\textbf{Value} & \textbf{Description} \\
\hline
$Q_5 = 226$ & 5th convergent denominator of $\pi/10$ \\
$a_6 = 733$ & 6th partial quotient of $\pi/10$ \\
$a_7 = 11$ & 7th partial quotient of $\pi/10$ \\
$82{,}829$ & Exact integer arithmetic: $733 \times 226 / 2$ \\
\hline
\end{tabular}
\end{center}

\noindent\textit{Values $Q_5 = 1296$ and $a_7 = 1$ are false and shall not appear in any certificate.}

\section*{3. Source Code}

\textbf{File:} \texttt{cf\_pi10.py}
\begin{lstlisting}
# Battle Plan v1.6 - Module 3: CF of pi/10
# Library: mpmath 1.3.0 (pinned for Modules 1-3)
# Seed fix: p=1,pp=0 tracks numerators; q=0,qq=1 tracks denominators
from mpmath import mp
mp.dps = 50

x = mp.pi/10
a = []
for _ in range(8):
    ai = int(x)
    a.append(ai)
    x = x - ai
    if x == 0: break
    x = 1/x

# a = [0, 3, 5, 2, 5, 1, 733, 11]
p, q = 1, 0
pp, qq = 0, 1
for i, ai in enumerate(a):
    p, pp = ai*p + pp, p
    q, qq = ai*q + qq, q
    if i == 5:  # n=5
        p5, q5 = p, q  # 71, 226

a6 = a[6]  # 733
a7 = a[7]  # 11
bound = a6 * q5 // 2  # 82829

print(f"CF: {a}")
print(f"p_5={p5}, Q_5={q5}, a_6={a6}, a_7={a7}, bound={bound}")
\end{lstlisting}

\noindent\textbf{Seed fix note:} The original LaTeX draft had seeds \texttt{p=0,pp=1,q=1,qq=0},
which swaps numerators and denominators, giving \texttt{p5=226} (wrong) and \texttt{q5=71} (wrong).
Fixed to \texttt{p=1,pp=0,q=0,qq=1}.

\section*{4. Build Environment}

\begin{lstlisting}
$ python3 --version
Python 3.10.12

$ pip show mpmath | grep Version
Version: 1.3.0

$ sha256sum cf_pi10.py
5ac750a4013027495d76ddd22fc76842f408ef716ba1498437d29414da8169b2  cf_pi10.py
\end{lstlisting}

\section*{5. Raw Execution Log}

\begin{lstlisting}
$ python3 cf_pi10.py
CF: [0, 3, 5, 2, 5, 1, 733, 11]
p_5=71, Q_5=226, a_6=733, a_7=11, bound=82829
\end{lstlisting}

\section*{6. Cryptographic Binding}

\begin{center}
\begin{tabular}{|l|l|}
\hline
\textbf{Item} & \textbf{SHA-256 Digest} \\
\hline
Source \texttt{cf\_pi10.py} & \texttt{5ac750a4013027495d76ddd22fc76842} \\
 & \texttt{f408ef716ba1498437d29414da8169b2} \\
\hline
Stdout (2 lines) & \texttt{e687bb09a55e4eda198d4c5b24d03b75} \\
 & \texttt{79f93bba27184a61fec7cbe29a83d044} \\
\hline
\end{tabular}
\end{center}

\section*{7. Causal Dependency}

Module 4 requires the inequality $p_5 > 82{,}829$ to prove that the enumeration of
$S(\alpha_0)$ is complete for $p \leq 10^{4000}$.
If SHA \texttt{e687bb09\ldots} changes, Module 4 must be re-verified.
The bound 474,984 is deprecated and invalid.

Module 3 has \emph{no} causal parent SHA --- it depends only on the mathematical
constant $\pi$ (computed via \texttt{mpmath} 1.3.0) and pure integer arithmetic.

\vspace{1em}
\noindent\textbf{Verification command:}
\begin{lstlisting}
$ python3 cf_pi10.py | sha256sum
e687bb09a55e4eda198d4c5b24d03b7579f93bba27184a61fec7cbe29a83d044  -
\end{lstlisting}

\vspace{0.5em}
\noindent\small Repository: \texttt{https://github.com/DavidFox998/alpha0-ponti} ·
Certificate generated by Machine Certificate v1.6

\end{document}
